\subsection{Independent ideal hybrid-MZI reference: do not set the amplitudes equal}
This comparison assumes an ideal PBS, two 50/50 NPBSs, no arm attenuation or
retardance, and zero phase-plate offsets. It is a reference limit, not an
assumption imposed on the configured network above. Let
$\alpha=\phi_1+\delta$ and retain independent input amplitudes $a_x,a_y$.
Propagation through the two combiners gives
\begin{align}
\mathbf E_E &= i e^{i\phi_2/2}
\begin{bmatrix}a_xe^{i\alpha}\sin(\phi_2/2)\\a_y\cos(\phi_2/2)\end{bmatrix},\\
\mathbf E_F &= i e^{i\phi_2/2}
\begin{bmatrix}a_xe^{i\alpha}\cos(\phi_2/2)\\-a_y\sin(\phi_2/2)\end{bmatrix}.
\end{align}
The shared phase has unit modulus and cancels in each field norm. Therefore
\begin{align}
I_E &= a_x^2\sin^2\!\frac{\phi_2}{2}+a_y^2\cos^2\!\frac{\phi_2}{2}
=\frac{a_x^2+a_y^2}{2}+\frac{a_y^2-a_x^2}{2}\cos\phi_2,\\
I_F &= a_x^2\cos^2\!\frac{\phi_2}{2}+a_y^2\sin^2\!\frac{\phi_2}{2}
=\frac{a_x^2+a_y^2}{2}-\frac{a_y^2-a_x^2}{2}\cos\phi_2,\\
I_E+I_F&=a_x^2+a_y^2,\qquad
\mathcal V_E=\mathcal V_F=\frac{|a_x^2-a_y^2|}{a_x^2+a_y^2}.
\end{align}
The visibility formula assumes a nonzero input and a full phase cycle.
Constant individual-port power follows only in the special case $a_x=a_y$.
Unequal effective amplitudes immediately produce complementary fringes, even
with ideal optical elements. Ideal PBS-to-NPBS1 arm transmissions $\ell_A,\ell_B$
can be included here by replacing $a_x$ with $\ell_A a_x$ and $a_y$ with
$\ell_B a_y$; their common attenuation and their ratio play different roles.

For the same independent amplitudes, the normalized Stokes coordinates are
\begin{align}
s_{1,E}&=\frac{a_y^2\cos^2(\phi_2/2)-a_x^2\sin^2(\phi_2/2)}{I_E},\\
s_{2,E}&=\frac{a_xa_y\sin\phi_2\cos\alpha}{I_E},\qquad
s_{3,E}=\frac{a_xa_y\sin\phi_2\sin\alpha}{I_E},\\
s_{1,F}&=\frac{a_y^2\sin^2(\phi_2/2)-a_x^2\cos^2(\phi_2/2)}{I_F},\\
s_{2,F}&=-\frac{a_xa_y\sin\phi_2\cos\alpha}{I_F},\qquad
s_{3,F}=-\frac{a_xa_y\sin\phi_2\sin\alpha}{I_F}.
\end{align}
These expressions are undefined at a dark output. Only after choosing
$a_x=a_y>0$ does port E reduce to
\begin{equation}
\mathbf s_E=(\cos\phi_2,\ \sin\phi_2\cos\alpha,\ \sin\phi_2\sin\alpha)^T.
\end{equation}
At fixed $\phi_2$, a $\phi_1$ scan has component amplitudes
$|\sin\phi_2|$ for both $s_2$ and $s_3$. At fixed $\phi_1$, a $\phi_2$ scan
has amplitudes $|\cos\alpha|$ and $|\sin\alpha|$ instead. Different heights
in these two cuts are expected; they do not indicate a loss of polarization.
The vector remains on the unit sphere.

\subsection{Why differential polarization phase can modulate total port power}
For a scalar final NPBS, write its incident fields as
$e^{i\phi_2}\mathbf c$ and $\mathbf d$, with real mixing coefficients
$t=\cos\mu$, $r=\sin\mu$ and $z=\mathbf c^\dagger\mathbf d$.
Direct expansion gives
\begin{align}
I_E&=t^2\|\mathbf c\|^2+r^2\|\mathbf d\|^2
+2tr[\operatorname{Re}z\sin\phi_2-\operatorname{Im}z\cos\phi_2],\\
I_F&=r^2\|\mathbf c\|^2+t^2\|\mathbf d\|^2
-2tr[\operatorname{Re}z\sin\phi_2-\operatorname{Im}z\cos\phi_2],\\
I_E+I_F&=\|\mathbf c\|^2+\|\mathbf d\|^2.
\end{align}
Thus the fringe is governed by the polarization overlap of the two incident
fields. Differential retardance can change that overlap without attenuating
either arm. Power moves between outputs while their sum remains fixed during
this phase scan.

As a transparent special case, keep the ideal PBS/NPBSs and add axis-aligned
retarders with retardances $\rho_C,\rho_D$ in arms C and D. Define
$\Delta\rho=\rho_C-\rho_D$. The E-port intensity becomes
\begin{equation}
I_E=\frac{a_x^2+a_y^2}{2}
+\frac{a_y^2\cos(\phi_2+\Delta\rho/2)
-a_x^2\cos(\phi_2-\Delta\rho/2)}{2}.
\end{equation}
Even in the explicitly chosen balanced case $a_x=a_y=a$, this gives
\begin{equation}
I_E=a^2[1-\sin\phi_2\sin(\Delta\rho/2)],\qquad
\mathcal V=|\sin(\Delta\rho/2)|.
\end{equation}
This aligned-retarder example explains the mechanism; rotated axes and other
imperfections must use the configured transfer calculation, not this special-case
formula. A fitted retarder parameter is not, by itself, an independently
measured property of an optic.
